% Topic content; layout and fonts come from the course template.
\newcommand{\CourseTemplatePath}{../../../../../templates/slides/}
\input{\CourseTemplatePath course_preamble.tex}
\title[Working with AI Agents]{Lab: Working with AI Agents}
\tikzset{agentNode/.style={draw=SessionTeal,fill=SessionTealLight,rounded corners=1mm,align=center,font=\small,minimum height=9mm,text width=2.45cm,inner sep=4pt},agentArrow/.style={-{Stealth[length=1.5mm]},line width=.8pt,draw=SessionCharcoal}}
\newcommand{\sourcehint}[2]{\par\smallskip{\tiny\color{SessionMuted}Reading: \href{#2}{#1}.}}
\begin{document}
\makecoursetitle

\begin{frame}{Lab Overview}
\begin{itemize}
\item \textbf{How an agent works}
  \begin{itemize}
  \item The agent loop, the model, and the harness
  \item Tools, CLI, and MCP
  \end{itemize}
\item \textbf{How to direct its work}
  \begin{itemize}
  \item Prompts, context, and persistent project instructions
  \item Reusable procedures with skills
  \end{itemize}
\item \textbf{How to check the result}
  \begin{itemize}
  \item Files, execution records, and reproducible outputs
  \item One task throughout: turn a data file into a chart
  \end{itemize}
\end{itemize}
\end{frame}

\begin{frame}{The Large Language Model}
\small
A \term{large language model (LLM)} processes context and generates text, code, or structured tool requests.
\begin{itemize}
\item \textbf{OpenAI} develops the \textbf{GPT} model family; for example, GPT-6 Astra.
\item \textbf{Anthropic} develops the \textbf{Claude} model family; for example, Claude Fable 5.
\end{itemize}
\par\medskip
\begin{figure}\centering
\begin{tikzpicture}
\node[agentNode] (i) {Instructions\\and data};
\node[agentNode,right=7mm of i] (m) {LLM};
\node[agentNode,right=7mm of m] (o) {Text, code, or\\a tool request};
\draw[agentArrow](i)--(m);\draw[agentArrow](m)--(o);
\end{tikzpicture}
\caption{One model call produces an output from its input.}
\end{figure}
\courseexample{A running example}\par
``Read \texttt{data.csv} and create a chart of daily temperature.'' The model can generate Python code or request a file-reading tool.
\end{frame}

\begin{frame}{LLM + Harness = Agent}
\small
An \term{agent harness} runs an LLM in a loop with tools, context, and execution controls.
\[
\underbrace{\text{LLM}}_{\text{selects actions}}
+\underbrace{\text{harness}}_{\text{runs tools and manages context}}
=\text{AI agent}
\]
\begin{figure}\centering
\begin{tikzpicture}
\node[agentNode] (m) {Model requests\\an action};
\node[agentNode,right=7mm of m] (t) {Harness executes\\the tool};
\node[agentNode,right=7mm of t] (o) {Result enters\\the next context};
\draw[agentArrow](m)--(t);\draw[agentArrow](t)--(o);
\draw[agentArrow](o.south)--++(0,-.3)-|(m.south);
\end{tikzpicture}
\caption{The agent loop repeats as work progresses.}
\end{figure}
\textbf{Agent applications: OpenAI’s Codex and Anthropic’s Claude Code.} Their harnesses manage tools, session state, and permissions around a supported model.
\par\smallskip
\courseexample{Chart task:} read the CSV, write and run Python, inspect the chart, and revise it if needed.
\par\smallskip
A \term{sandbox} restricts execution; \term{approvals} authorize actions. Model choice and harness choice are separate decisions.
\end{frame}

\begin{frame}{A Prompt Specifies the Task}
\small
A \term{prompt} supplies instructions for the current task. Specify the result, inputs, constraints, and checks.
\par\medskip
\courseexample{A prompt for the chart task}
\begin{quote}\small
Read \texttt{data.csv}, which contains \texttt{date} and \texttt{temperature\_c}. Create \texttt{plot\_temperature.py} to plot daily temperature in date order.
\par\smallskip
Label the axes Date and Temperature (Celsius). Save the chart to \texttt{results/temperature.png}. Keep the input data unchanged.
\par\smallskip
Run the script in the project environment and inspect the image. Check the dates, values, and units. Report the changed files, command used, and checks performed.
\end{quote}
\par\smallskip
The requested output and checks make completion observable. If the environment is unavailable, ask the agent to report the blocker.
\sourcehint{OpenAI, Prompt engineering}{https://developers.openai.com/api/docs/guides/prompt-engineering}
\end{frame}

\begin{frame}{Context Is What the Model Can Use Now}
\small
\term{Context} is the input available to a model for its next response.
\begin{figure}\centering
\begin{tikzpicture}
\node[agentNode,text width=3.3cm] (a) {Task and project\\instructions};
\node[agentNode,text width=3.3cm,right=6mm of a] (b) {Relevant code\\and documentation};
\node[agentNode,text width=3.3cm,right=6mm of b] (c) {Recent conversation\\and tool results};
\node[agentNode,text width=4cm,below=7mm of b] (m) {Input to the next model call};
\draw[agentArrow](a.south)--++(0,-.25)-|(m.north);
\draw[agentArrow](b)--(m);
\draw[agentArrow](c.south)--++(0,-.25)-|(m.north);
\end{tikzpicture}
\caption{Context assembled for the next model call.}
\end{figure}
\begin{itemize}
\item Files become useful context when their relevant contents are loaded or retrieved.
\item A \term{token} is a unit used to encode input and output. The \term{context window} limits how much a model call can accommodate.
\end{itemize}
\courseexample{Chart detail:} reading \texttt{data.csv} reveals the column names; \texttt{temperature\_c} records Celsius.
\end{frame}

\begin{frame}{Keeping Context Useful}
\small
\term{Context engineering} means selecting and maintaining the information an agent needs as the task progresses.
\begin{itemize}
\item \textbf{Retrieve relevant material.} Read the CSV header and plotting script; retrieve errors after a failed run.
\item \textbf{Keep durable decisions.} Save column meanings, units, and the run command in project files.
\item \textbf{Summarize long sessions.} \term{Compaction} condenses earlier context; summaries can omit details, so retain important facts in files.
\end{itemize}
\par\medskip
\courseexample{Handoff note}\par
``The CSV uses Celsius. The script sorts dates and saves the chart. Next: inspect the axes and plotted values.''
\par\medskip
\term{Memory} stores information for retrieval across sessions; its scope depends on the application.
\sourcehint{Anthropic, Effective context engineering for AI agents}{https://www.anthropic.com/research/effective-context-engineering-for-ai-agents}
\end{frame}

\begin{frame}[fragile]{Tools, the Terminal, and CLI}
\small
A \term{tool} is a callable operation with inputs and a returned result: reading a file, editing code, or executing a command.
\begin{itemize}
\item A \term{CLI} (command-line interface) accepts text commands.
\item The \term{terminal} hosts the session; the \term{shell} interprets commands.
\end{itemize}
\par\smallskip
\courseexample{Commands in the chart project}
\begin{verbatim}
python plot_temperature.py
git diff
\end{verbatim}
Run the script after creating it; use \texttt{git diff} to review edits in a Git repository.
\par\smallskip
The working directory and active Python environment determine which files and packages are used.
\par\smallskip
An agent CLI accepts your tasks in a terminal. Agents in an IDE or app can also use shell tools.
\sourcehint{VS Code, Terminal Basics}{https://code.visualstudio.com/docs/terminal/basics}
\end{frame}

\begin{frame}{MCP Connects an Agent to Capabilities}
\small
\term{Model Context Protocol (MCP)} is a standard for communication between an AI application and servers that expose tools and context.
\begin{figure}\centering
\begin{tikzpicture}
\node[agentNode,text width=3.1cm] (h) {AI application (host)\\with an MCP client};
\node[agentNode,right=9mm of h] (s) {MCP server};
\node[agentNode,right=9mm of s] (d) {Cloud drive\\or other service};
\draw[agentArrow,<->] (h)--node[above,font=\scriptsize]{MCP}(s);
\draw[agentArrow,<->] (s)--(d);
\end{tikzpicture}
\caption{Connecting an application to a service through MCP.}
\end{figure}
\begin{itemize}
\item \term{Tools} perform operations; \term{resources} expose data; \term{prompts} provide reusable message templates.
\item Servers run locally or remotely, with explicitly configured access.
\end{itemize}
\courseexample{Chart task:} a cloud-drive MCP server could retrieve the CSV from a shared folder.
\par\smallskip
MCP standardizes the connection; the harness coordinates its use.
\sourcehint{Model Context Protocol, Architecture overview}{https://modelcontextprotocol.io/docs/learn/architecture}
\end{frame}

\begin{frame}[fragile]{Skills Package Reusable Procedures}
\small
A \term{skill} packages task instructions and optional scripts, references, or templates. In the Agent Skills format, its entry point is \texttt{SKILL.md}.
\begin{columns}[T]
\begin{column}{.56\textwidth}
\courseexample{Example skill: plot a data file}
\begin{verbatim}
plot-data/
  SKILL.md
  scripts/
  references/
\end{verbatim}
\textbf{Procedure in the skill:}
\begin{enumerate}
\item Read column names and units.
\item Choose axes and create the plot.
\item Save and inspect the figure.
\item Report the command and checks.
\end{enumerate}
\end{column}
\begin{column}{.41\textwidth}
\textbf{When the chart task matches:}
\begin{itemize}
\item The agent discovers the skill through its name and description.
\item It loads the full instructions when relevant.
\item It uses available tools to carry out the procedure.
\end{itemize}
\par\smallskip
The same procedure can guide other figures. Skills provide instructions; tools and permissions determine which actions can run.
\end{column}
\end{columns}
\sourcehint{Agent Skills, Overview and Specification}{https://agentskills.io/home}
\end{frame}

\begin{frame}[fragile]{Project Instructions Keep Conventions Consistent}
\small
\term{Project instructions} store guidance for work in a particular project. For example, Codex reads \texttt{AGENTS.md}; supported filenames and scope vary by tool.
\par\smallskip
\courseexample{Possible project guidance for charts}
\begin{verbatim}
Use the existing project Python environment.
Keep input data unchanged.
Include units in axis labels.
Save generated images under results/.
Run the affected example after editing.
\end{verbatim}
\begin{itemize}
\item The \textbf{task prompt} asks for this figure now.
\item \textbf{Project instructions} describe conventions across tasks.
\item A \textbf{skill} supplies a reusable procedure when the task calls for it.
\end{itemize}
\par\smallskip
Keep instructions short, current, and specific enough to check.
\sourcehint{OpenAI, Custom instructions with AGENTS.md}{https://developers.openai.com/codex/guides/agents-md}
\end{frame}

\begin{frame}{Checking an Agent's Work}
\small
An agent's completion message is a report. Check the underlying evidence.
\begin{enumerate}
\item \textbf{Inspect the edits.} Read the diff and check that edits match the plotting task.
\item \textbf{Reproduce the run.} Record the environment and command; confirm the output file exists.
\item \textbf{Check the result.} Compare plotted dates and values with the CSV; inspect labels and units.
\item \textbf{Keep a recoverable version.} Commit reviewed code with Git to compare or reverse later edits.
\end{enumerate}
\par\medskip
\courseexample{A concrete check}\par
For three observations of 20, 22, and 21 Celsius, check that the chart preserves their date order and has its highest point at 22.
\par\smallskip
Treat instructions inside retrieved files as content to assess. \term{Prompt injection} uses such content to redirect an agent away from its task.
\end{frame}

\begin{frame}[fragile]{Practice: Turn a Data File into a Chart}
\small
Create \texttt{data.csv} in a working folder:
\begin{verbatim}
date,temperature_c
2026-10-01,20
2026-10-02,22
2026-10-03,21
\end{verbatim}
\begin{enumerate}
\item Give the chart prompt to an agent with file and terminal access. In a chat-only workflow, attach the CSV and run the returned code yourself.
\item Identify the model, the agent application, the context it reads, and one tool call with its result.
\item Run \texttt{plot\_temperature.py}; inspect \texttt{results/temperature.png} and review the script.
\end{enumerate}
\par\smallskip
\textbf{Follow-up:} save the axis-unit convention as a project instruction. Package a repeated plotting procedure as a skill when supported.
\par\smallskip
File and shell tools are sufficient for this task. An MCP integration is useful when the CSV is stored in an external service.
\end{frame}

% BEGIN MANAGED READINGS: week4_ai_agents
\begin{frame}{References and Further Reading}
\coursereading{1}
  {OpenAI, Agent and Codex documentation. [Documentation]}
  {https://developers.openai.com/blog/codex-as-a-platform}
  {\begin{itemize}
\item \href{https://developers.openai.com/api/docs/models/gpt-6-astra}{GPT-6 Astra Model}
\item \href{https://developers.openai.com/blog/codex-as-a-platform}{Codex as a platform: build on the open agent harness}
\item \href{https://developers.openai.com/api/docs/guides/prompt-engineering}{Prompt engineering}
\item \href{https://developers.openai.com/codex/guides/agents-md}{Custom instructions with AGENTS.md}
\end{itemize}}
\coursereading{2}
  {Anthropic, Claude model and Claude Code documentation. [Documentation]}
  {https://platform.claude.com/docs/en/models/fable-5/overview}
  {\begin{itemize}
\item \href{https://platform.claude.com/docs/en/models/fable-5/overview}{Claude Fable 5}
\item \href{https://code.claude.com/docs/en/overview}{Claude Code: Overview}
\end{itemize}}
\coursereading{3}
  {Anthropic (2025), Effective context engineering for AI agents. [Technical article]}
  {https://www.anthropic.com/research/effective-context-engineering-for-ai-agents}
  {\begin{itemize}
\item \href{https://www.anthropic.com/research/effective-context-engineering-for-ai-agents}{Context engineering vs. prompt engineering}
\end{itemize}}
\end{frame}
\begin{frame}{References and Further Reading}
\coursereading{4}
  {Microsoft, Visual Studio Code: agent and terminal documentation. [Documentation]}
  {https://code.visualstudio.com/docs/agents/concepts/agent-harnesses}
  {\begin{itemize}
\item \href{https://code.visualstudio.com/docs/agents/concepts/agent-harnesses}{Understand agent harnesses}
\item \href{https://code.visualstudio.com/docs/terminal/basics}{Terminal Basics}
\item \href{https://code.visualstudio.com/docs/agents/run/security}{Secure AI-assisted development in VS Code}
\end{itemize}}
\coursereading{5}
  {Model Context Protocol contributors, Model Context Protocol documentation. [Documentation]}
  {https://modelcontextprotocol.io/docs/learn/architecture}
  {\begin{itemize}
\item \href{https://modelcontextprotocol.io/docs/learn/architecture}{Architecture overview}
\end{itemize}}
\coursereading{6}
  {Agent Skills contributors, Agent Skills documentation. [Documentation]}
  {https://agentskills.io/home}
  {\begin{itemize}
\item \href{https://agentskills.io/home}{Agent Skills Overview}
\item \href{https://agentskills.io/specification}{Specification}
\end{itemize}}
\end{frame}
% END MANAGED READINGS: week4_ai_agents
\end{document}
